% =====================================================================
% P1 Appendix A.3 --- Theorem 2 (finite-sample coverage under drift).
%
% Source of truth: paper/theorem2.md v0.5 (2026-10-21 theorem statements/proofs; 2026-10-26 interface wording).
% This file is a fragment: it is meant to be \input under
%   \appendix \section{Proofs}
% after A.1 and A.2 (the standalone driver main-a3.tex sets the
% subsection counter to 2, so this fragment becomes A.3).
%
% Requirements: amsmath, amssymb, amsthm, natbib, and url (hyperref loads
% url). Theorem-like environments are defined here with the explicit paper
% numbering (Theorem 2A/2B, Corollaries 2.1/2.2, Lemmas 1--3,
% Remarks 2.1--2.3 and 3.1). Equation tags (2.1)--(2.6) and (3.1)--(3.3)
% follow the frozen source; do not renumber without updating the main text.
% =====================================================================

\newtheorem*{aIIIthmA}{Theorem 2A}
\newtheorem*{aIIIthmB}{Theorem 2B}
\newtheorem*{aIIIcorA}{Corollary 2.1}
\newtheorem*{aIIIcorB}{Corollary 2.2}
\newtheorem*{aIIIlemA}{Lemma 1}
\newtheorem*{aIIIlemB}{Lemma 2}
\newtheorem*{aIIIlemC}{Lemma 3}
\newtheorem*{aIIIremA}{Remark 2.1}
\newtheorem*{aIIIremB}{Remark 2.2}
\newtheorem*{aIIIremC}{Remark 2.3}
\newtheorem*{aIIIremD}{Remark 3.1}

% Fixed equation tags do not advance the equation counter, so hyperref would
% give every display the same destination. A dedicated counter keeps the
% destinations unique without changing the printed (frozen) tag numbers. The
% original \theHequation is restored at the end of the fragment.
\providecommand{\theHequation}{}
\let\aIIItheHequation\theHequation
\newcounter{aIIIanchor}
\renewcommand{\theHequation}{aIII.\arabic{aIIIanchor}}
\newcommand{\aIIIfixedtag}[1]{\stepcounter{aIIIanchor}\tag{#1}}

\subsection{Finite-sample coverage under drift (Theorem~2)}\label{app:a3}

This appendix gives the finite-sample statement and proof of Theorem~2 in the
notation of the main text. It is \emph{not} a population
identification result. Theorem~2A treats marginal coverage of split conformal when
the drift $\kappa$ is known; Theorem~2B treats conditional (PAC) coverage when
$\kappa$ is estimated; Corollary~2.2 gives a confidence lower bound on the marginal
coverage rate using only $\hat\kappa$; and Remark~2.3 gives a fallback bound in
expectation. None of these results depends on the POMDP identification of
Theorem~1: when the scores are produced by the plug-in pipeline of Theorem~1, all
statements are conditional on the fixed score function.

\paragraph{Notation.}
Let $E_1,\dots,E_n$ be i.i.d.\ with CDF $F_{\mathrm{cal}}$, let
$E'_1,\dots,E'_m$ be i.i.d.\ with CDF $F_{\mathrm{test}}$, and let
$E\sim F_{\mathrm{test}}$ be independent of the calibration sample. Write
$\hat F_{\mathrm{cal}}$ and $\hat F_{\mathrm{test}}$ for the empirical CDFs, and
write $\kappa:=\kappa_{\mathrm{conf}}$ and $\hat\kappa:=\hat\kappa_{\mathrm{conf}}$
for the drift and its estimate used in the main text:
\[
\kappa=\sup_x\bigl|F_{\mathrm{cal}}(x)-F_{\mathrm{test}}(x)\bigr|,
\qquad
\hat\kappa=\sup_x\bigl|\hat F_{\mathrm{cal}}(x)-\hat F_{\mathrm{test}}(x)\bigr|,
\]
where $\hat\kappa$ is the two-sample Kolmogorov--Smirnov statistic.

Fix $\alpha\in(0,1)$ and set $k=\lceil(n+1)(1-\alpha)\rceil$. If $k\le n$, let
$\hat q:=E_{(k)}$ be the $k$-th order statistic of the calibration scores; if
$k>n$, set $\hat q:=+\infty$. The prediction set is
$C(x)=\{y:s(x,y)\le\hat q\}$ and $E=s(X,Y)$ is the conformity score; only the
marginal law of the scores is used and the specific form of $s$ is irrelevant.
Raw, studentized and CQR scores are special cases. Let
$\delta_{\mathrm{cal}},\delta_\kappa\in(0,1)$ be confidence parameters, and let
\[
\gamma:=\Pr\{E\le\hat q\}
\]
be the \emph{marginal coverage rate} of the split-conformal procedure, where the
probability is over the calibration sample and the new test point. Given
$F_{\mathrm{cal}},F_{\mathrm{test}},n,k$, $\gamma$ is a constant; it does not vary
with a particular calibration draw.

\paragraph{Reading convention.}
The coverage probability in Theorem~2A is taken over the calibration sample and the
new test point (marginal coverage), i.e.\ it is $\gamma$; it is \emph{not} a
pathwise lower bound on $F_{\mathrm{cal}}(\hat q)$ for every calibration sample.
The quantity $F_{\mathrm{test}}(\hat q)$ in Theorem~2B is the conditional coverage
given the data, and its $1-\delta$ lower bound holds with probability over the
calibration and test-score samples; it does \emph{not} imply an unconditional bound
of the same value on $\Pr\{E\le\hat q\}$. Corollary~2.2 combines the constant bound
of 2A with a confidence upper bound on $\kappa$ and yields a $1-\delta_\kappa$ lower
confidence bound on $\gamma$; Remark~2.3 gives an unconditional bound in
expectation.

\begin{aIIIthmA}
Let $E_1,\dots,E_n$ be i.i.d.\ $F_{\mathrm{cal}}$ and let $E$ be independent of the
calibration sample with $E\sim F_{\mathrm{test}}$. Set
$\kappa=\sup_x|F_{\mathrm{cal}}(x)-F_{\mathrm{test}}(x)|$. Then for
$k=\lceil(n+1)(1-\alpha)\rceil\le n$,
\begin{equation}\label{eq:a3-2.1}
\Pr\{E\le\hat q\}\;\ge\;\max\!\Bigl(0,\;\frac{k}{n+1}-\kappa\Bigr)
\;\ge\;\max\!\bigl(0,\;1-\alpha-\kappa\bigr).
\aIIIfixedtag{2.1}
\end{equation}
If $k>n$, then $\hat q=+\infty$ and the coverage is $1$.
\end{aIIIthmA}

The calibration and test distributions need not be exchangeable; only the i.i.d.\
calibration sample and the independence of the test point are required. The term
$k/(n+1)$ is the standard finite-sample term of split conformal
\cite{vovk2005,lei2018} (with equality in the continuous case), and $-\kappa$ is
the price of drift. This is the exact form of Proposition~7; no ``ignoring
finite-sample terms'' is involved. If only an upper bound
$\bar\kappa\ge\kappa$ is known (the original $\kappa_{\mathrm{conf}}$ formulation),
replace $\kappa$ by $\bar\kappa$ in \eqref{eq:a3-2.1}.

\begin{aIIIremA}
The proof only uses
$|F_{\mathrm{test}}(\hat q)-F_{\mathrm{cal}}(\hat q)|\le\kappa$ at the random
threshold $\hat q$; the supremum over $x$ is taken so that $\kappa$ is estimable by
a two-sample KS statistic. More precisely, 2A only needs the one-sided inequality
$F_{\mathrm{cal}}(\hat q)-F_{\mathrm{test}}(\hat q)\le\kappa$. When the direction of
drift is known, the sharper one-sided quantity
$\kappa^+:=\sup_x(F_{\mathrm{cal}}(x)-F_{\mathrm{test}}(x))$ can replace $\kappa$;
its estimated version requires a one-sided DKW inequality and is not developed
here. A local CDF difference near $\hat q$ would give a smaller quantity, but the
result would no longer be a distribution-free KS bound.
\end{aIIIremA}

\begin{aIIIthmB}
In the setting of 2A, suppose in addition that $E'_1,\dots,E'_m$ are i.i.d.\
$F_{\mathrm{test}}$, independent of the calibration sample and of $E$, and let
$\hat\kappa$ be the two-sample KS statistic. For any
$\delta_{\mathrm{cal}},\delta_\kappa\in(0,1)$, with probability at least
$1-\delta_{\mathrm{cal}}-\delta_\kappa$ (over the calibration and test-score
samples),
\begin{equation}\label{eq:a3-2.2}
F_{\mathrm{test}}(\hat q)\;\ge\;
\max\!\Bigl(0,\;\frac{k}{n}-\varepsilon_{\mathrm{cal}}-\hat\kappa-\varepsilon_\kappa\Bigr),
\aIIIfixedtag{2.2}
\end{equation}
where
\begin{equation}\label{eq:a3-2.3}
\varepsilon_{\mathrm{cal}}=\sqrt{\frac{\ln(2/\delta_{\mathrm{cal}})}{2n}},
\qquad
\varepsilon_\kappa=\sqrt{\frac{\ln(4/\delta_\kappa)}{2n}}
+\sqrt{\frac{\ln(4/\delta_\kappa)}{2m}}.
\aIIIfixedtag{2.3}
\end{equation}
\end{aIIIthmB}

This is a PAC bound on \emph{conditional} coverage, not an unconditional marginal
bound; for the PAC/risk-control reading of such statements see
\cite{bates2021}. Equation~\eqref{eq:a3-2.2} says that, with probability at least
$1-\delta_{\mathrm{cal}}-\delta_\kappa$ over the calibration and test-score samples,
the coverage probability
$F_{\mathrm{test}}(\hat q)=\Pr\{E\le\hat q\mid\text{data}\}$ of a new test point is
at least the right-hand side. It does not imply that $\Pr\{E\le\hat q\}$, with the
probability also over the data, is bounded below by the same value. To see the
distinction, write the right-hand side as the random variable $b$: on the good event
$A\cap B$ we have $F_{\mathrm{test}}(\hat q)\ge b$, and on the bad event
$F_{\mathrm{test}}(\hat q)\ge0$, so the unconditional coverage is at least
$\mathbb E[b\,\mathbf 1_{A\cap B}]$. Since $b\le1$, its worst case (with the bad
event concentrated where $b$ is largest) is $\mathbb E[b]-\delta$; when $b$ is
constant this reduces to $(1-\delta)b$. Note that $(1-\delta)\mathbb E[b]$ is
generally \emph{not} valid for random $b$: if $b\in\{0,1\}$ and
$\Pr\{b=1\}=p\le\delta$, the left-hand side can be $0$. The exact expectation bound
is Remark~2.3; for a marginal confidence bound that uses only $\hat\kappa$, see
Corollary~2.2.

\paragraph{Why $k/n$ rather than $k/(n+1)$.}
The term $k/(n+1)$ in 2A is an identity in \emph{expectation},
$\mathbb E[F_{\mathrm{cal}}(\hat q)]\ge k/(n+1)$ (Lemma~1). Theorem~2B needs a
pathwise lower bound given the data and can only use
$\hat F_{\mathrm{cal}}(\hat q)\ge k/n$ (the $k$-th order statistic is exceeded by at
most $n-k$ calibration scores) together with the calibration DKW slack
$\varepsilon_{\mathrm{cal}}$. The difference is $k/(n(n+1))\le1/n$, about $0.007$
at $n=128$.

\begin{aIIIremB}
If the test scores \emph{and their coverage outcomes} are observable, then
$\hat F_{\mathrm{test}}(\hat q)$ is the empirical coverage on the test sample and,
for any $\delta\in(0,1)$, the one-sample DKW inequality gives
\begin{equation}\label{eq:a3-2.4}
F_{\mathrm{test}}(\hat q)\;\ge\;
\hat F_{\mathrm{test}}(\hat q)-\sqrt{\frac{\ln(2/\delta)}{2m}}
\quad\text{with probability }1-\delta,
\aIIIfixedtag{2.4}
\end{equation}
where the probability is over the $m$ test scores and $\hat q$ is held fixed
(calibration and test are independent).
This bound is typically sharper than \eqref{eq:a3-2.2} (about $0.75$ for the
studentized score in our experiment, versus $0.38$ for \eqref{eq:a3-2.2}), but it
uses only the coverage on the test sample itself and does not transfer to a
deployment setting without test labels. The value of $\kappa$ is that it is a
transferable drift quantity. The main text reports the PAC bound
\eqref{eq:a3-2.2} and the $\hat\kappa$-based marginal confidence bound of
Corollary~2.2; the direct test-CDF bound is stated here and is not part of the
main-text table.
\end{aIIIremB}

\begin{aIIIremC}
With the notation of 2B, let
$b:=\max(0,k/n-\varepsilon_{\mathrm{cal}}-\hat\kappa-\varepsilon_\kappa)$ and
$\delta:=\delta_{\mathrm{cal}}+\delta_\kappa$. Then
\begin{equation}\label{eq:a3-2.5}
\Pr\{E\le\hat q\}\;\ge\;\max\!\bigl(0,\;\mathbb E[b]-\delta\bigr).
\aIIIfixedtag{2.5}
\end{equation}
Indeed, on the good event $A\cap B$ (probability at least $1-\delta$) we have
$F_{\mathrm{test}}(\hat q)\ge b$, while on the bad event
$F_{\mathrm{test}}(\hat q)\ge0$; hence pathwise
$F_{\mathrm{test}}(\hat q)\ge b\,\mathbf 1_{A\cap B}$. Taking expectations and using
$b\le1$,
$\mathbb E[F_{\mathrm{test}}(\hat q)]\ge\mathbb E[b]-\Pr\{(A\cap B)^c\}\ge
\mathbb E[b]-\delta$. The bound involves $\mathbb E[b]$, not the observed $b$; in
practice one estimates $\mathbb E[b]$ by bootstrap and reports a CI, but this is
\emph{heuristic} ($\hat\kappa$ is a sup-type statistic and the bootstrap is biased
for non-pivotal functionals). The bound is usually weaker than
\eqref{eq:a3-2.2} and Corollary~2.2; the main text therefore states
\eqref{eq:a3-2.2} as the conditional PAC bound and keeps \eqref{eq:a3-2.5} in the
appendix as a fallback unconditional statement in expectation. For a marginal
confidence bound that uses only $\hat\kappa$, use Corollary~2.2.
\end{aIIIremC}

\begin{aIIIcorA}
Replacing $\kappa$ by $\hat\kappa$ and $k/(n+1)$ by $1-\alpha$ in 2A gives
$\max(0,1-\alpha-\hat\kappa)$. This is a \emph{plug-in point estimate}: it neither
accounts for the sampling error of $\hat\kappa$ nor retains the finite-sample term
$k/(n+1)$. Finite-sample \emph{conditional} coverage requires 2B; finite-sample
\emph{marginal} coverage requires 2A (if $\kappa$ is known), Corollary~2.2 (a
confidence bound with estimated $\kappa$), or the expectation bound of Remark~2.3.
Substituting $\hat\kappa$ directly does not give a guarantee.
\end{aIIIcorA}

This is precisely the risk of using an estimator as an upper bound. The quantity
$\max(0,1-\alpha-\hat\kappa)$ is an estimate whose uncertainty is
described by a bootstrap CI; the distribution-free conditional PAC lower bound is
\eqref{eq:a3-2.2}; the marginal coverage confidence lower bound that uses only
$\hat\kappa$ is Corollary~2.2; and the fallback in expectation is
\eqref{eq:a3-2.5} (weaker than Corollary~2.2 and requiring an estimate of
$\mathbb E[b]$). At $n=128$, \eqref{eq:a3-2.2} is empty for the raw and CQR scores,
whereas Corollary~2.2 still gives a small positive bound for CQR; see the main
text.

\begin{aIIIcorB}
In the setting of 2A, let $\gamma:=\Pr\{E\le\hat q\}$ be the marginal coverage rate
of the split-conformal procedure (probability over the calibration sample and an
independent new test point; given $F_{\mathrm{cal}},F_{\mathrm{test}},n,k$ it is a
constant). Let $E'_1,\dots,E'_m$ be test scores used for $\hat\kappa$, and let
$\varepsilon_\kappa$ be as in \eqref{eq:a3-2.3}. Then for any
$\delta_\kappa\in(0,1)$, with probability at least $1-\delta_\kappa$ over the
calibration and test-score samples,
\begin{equation}\label{eq:a3-2.6}
\gamma\;\ge\;
\max\!\Bigl(0,\;\frac{k}{n+1}-\hat\kappa-\varepsilon_\kappa\Bigr).
\aIIIfixedtag{2.6}
\end{equation}
\end{aIIIcorB}
\begin{proof}
By 2A, $\gamma\ge k/(n+1)-\kappa$. This is a deterministic inequality for the
\emph{constant} $\gamma$, not a pathwise lower bound on $F_{\mathrm{cal}}(\hat q)$
for an individual calibration sample. By Lemma~3, the event
$A=\{\kappa\le\hat\kappa+\varepsilon_\kappa\}$ has probability at least
$1-\delta_\kappa$. On $A$,
$\gamma\ge k/(n+1)-\kappa\ge k/(n+1)-\hat\kappa-\varepsilon_\kappa$.
\end{proof}

\paragraph{Reading notes.}
\begin{enumerate}
\item The object of \eqref{eq:a3-2.6} is the \emph{average coverage rate} $\gamma$
  of the procedure, not the coverage probability of the particular interval in
  hand. It answers ``what is the marginal coverage of this split-conformal
  procedure over repeated calibration draws, at confidence $1-\delta_\kappa$''
  and not ``given this batch of data, what is the conditional coverage of the next
  test point'' (the latter is 2B).
\item Equation~\eqref{eq:a3-2.1} must \emph{not} be read as the pathwise statement
  $F_{\mathrm{cal}}(\hat q)\ge k/(n+1)$. The term $k/(n+1)$ comes from the
  \emph{expectation} in Lemma~1, and $F_{\mathrm{cal}}(\hat q)$ can fall below it
  for an individual calibration sample. The proof of Corollary~2.2 applies 2A only
  to the constant $\gamma$ and makes no assertion about the realized
  $F_{\mathrm{cal}}(\hat q)$. This is the most common place to err.
\item Numerically, \eqref{eq:a3-2.6} is usually stronger than \eqref{eq:a3-2.2}
  (it omits $\varepsilon_{\mathrm{cal}}$, uses $k/(n+1)$ instead of $k/n$, and
  spends confidence only on $\delta_\kappa$), but it concerns a \emph{marginal}
  object and does not replace the conditional semantics of 2B. The two are
  reported side by side.
\item If $\kappa$ is known, \eqref{eq:a3-2.6} reduces to 2A. If $\hat\kappa$ comes
  from a sample that overlaps the evaluated test points (for example a pooled
  two-sample KS statistic), \eqref{eq:a3-2.6} remains a valid confidence bound:
  $\gamma$ is a population constant, so the overlap affects only the honesty of
  $m$, not validity. The strict version uses the fold's own $n_{\mathrm{cal}}$
  scores at $\tau=2$ and the fold's own $m$ scores at $\tau=81$; see the
  implementation notes below.
\item The proof only requires each sample to be i.i.d.\ \emph{within itself} (the
  two one-sample DKW inequalities are applied separately); independence between
  the two samples and of $E$ is a sufficient condition, not a necessary one (the
  union bound does not require independence).
\end{enumerate}

\paragraph{Relation to conditional coverage.}
Theorem~2A is \emph{marginal}: it averages over the test point according to
$F_{\mathrm{test}}$ and does not condition on $X$, regime or branch.
Corollary~2.2 is a confidence lower bound on the same marginal object $\gamma$
with estimated $\kappa$, and Theorem~2B is a PAC lower bound on the
\emph{conditional} coverage given the data. None of the three is a
branch-conditional coverage guarantee; see \cite{barber2021} for the limits
of distribution-free conditional inference. The drift $\kappa$ is a
one-dimensional KS distance between marginal score CDFs and is \emph{not} a
sufficient statistic for conditional coverage; the non-identifiability and minimax
lower bounds for branch-conditional coverage are in Theorem~3 and Proposition~2
(Propositions~8--12). In particular, \eqref{eq:a3-2.1}, \eqref{eq:a3-2.2} and
\eqref{eq:a3-2.6} must not be used to claim that every regime is covered.

\subsubsection{Auxiliary lemmas}

\begin{aIIIlemA}
Let $E_1,\dots,E_n$ be i.i.d.\ $F$, let $E^*$ be an independent copy, and let
$E_{(k)}$ be the $k$-th order statistic ($1\le k\le n$). Then
\begin{equation}\label{eq:a3-3.1}
\mathbb E\bigl[F(E_{(k)})\bigr]\;\ge\;\frac{k}{n+1}.
\aIIIfixedtag{3.1}
\end{equation}
\end{aIIIlemA}
\begin{proof}
Sort the $n+1$ i.i.d.\ variables $E^*,E_1,\dots,E_n$, breaking ties uniformly at
random, and let $R$ be the rank of $E^*$ (with $1$ the smallest). By
exchangeability, $R$ is uniform on $\{1,\dots,n+1\}$. If $R\le k$, then $E^*$ lies
among the first $k$ in the tie-broken order, so $E^*\le E_{(k)}$ (ties only make the
inequality easier). Therefore
\[
\mathbb E[F(E_{(k)})]
=\Pr\{E^*\le E_{(k)}\}
\ge\Pr\{R\le k\}
=\frac{k}{n+1}.
\]
In the continuous case $F(E_{(k)})\sim\mathrm{Beta}(k,n+1-k)$ and the expectation
equals $k/(n+1)$; with ties, random tie-breaking leaves the conclusion unchanged.
A discrete example shows the bound can be conservative: for a Bernoulli$(1/2)$
law with $n=k=1$, $\mathbb E[F(E_{(1)})]=3/4>1/2$.
\end{proof}

\begin{aIIIlemB}[\cite{massart1990}]
Let $E_1,\dots,E_N$ be i.i.d.\ $F$ and let $\hat F_N$ be the
empirical CDF. Then for every $\varepsilon>0$,
\begin{equation}\label{eq:a3-3.2}
\Pr\Bigl\{\sup_x|\hat F_N(x)-F(x)|>\varepsilon\Bigr\}
\le 2e^{-2N\varepsilon^2}.
\aIIIfixedtag{3.2}
\end{equation}
Equivalently, with probability at least $1-\delta$,
$\sup_x|\hat F_N-F|\le\sqrt{\ln(2/\delta)/(2N)}$.
\end{aIIIlemB}

\begin{aIIIlemC}
Under the independent-sample setting of 2B, for every
$\delta_\kappa\in(0,1)$,
\begin{equation}\label{eq:a3-3.3}
\Pr\{\kappa\le\hat\kappa+\varepsilon_\kappa\}\;\ge\;1-\delta_\kappa,
\qquad
\varepsilon_\kappa=
\sqrt{\frac{\ln(4/\delta_\kappa)}{2n}}
+\sqrt{\frac{\ln(4/\delta_\kappa)}{2m}}.
\aIIIfixedtag{3.3}
\end{equation}
\end{aIIIlemC}
\begin{proof}
By the triangle inequality,
\[
\sup_x\bigl|(F_{\mathrm{cal}}-F_{\mathrm{test}})
-(\hat F_{\mathrm{cal}}-\hat F_{\mathrm{test}})\bigr|
\le
\sup_x|\hat F_{\mathrm{cal}}-F_{\mathrm{cal}}|
+\sup_x|\hat F_{\mathrm{test}}-F_{\mathrm{test}}|.
\]
Apply Lemma~2 to the two terms with failure probability $\delta_\kappa/2$ each:
each term is at most $\sqrt{\ln(4/\delta_\kappa)/(2n)}$ and
$\sqrt{\ln(4/\delta_\kappa)/(2m)}$, respectively, with probability at least
$1-\delta_\kappa/2$; the union bound gives \eqref{eq:a3-3.3}.
\end{proof}

\begin{aIIIremD}
The frequently cited form $2\exp(-2nm\varepsilon^2/(n+m))$ is the \emph{two-sample}
DKWM inequality. \cite{weiDudley2012} show that, for equal sample sizes, the
constant $2$ is valid \emph{if and only if} $n\ge458$; there are counterexamples for
smaller samples. We therefore do not use that form and instead combine two
\emph{one-sample} DKW inequalities as in \eqref{eq:a3-3.3}, which is valid for all
$n,m$ at the price of a larger slack. At $n=m=128$ the price is $0.28$ versus
$0.19$, but it buys a genuine finite-sample guarantee; moreover
$n=m=128<458$ lies exactly in the uncovered regime of the constant-$2$
two-sample form.
\end{aIIIremD}

\subsubsection{Proof of Theorem~2A}

Given the calibration sample, $\hat q$ is fixed and $E$ is independent of the
calibration sample, so
\[
\Pr\{E\le\hat q\mid E_1,\dots,E_n\}=F_{\mathrm{test}}(\hat q)
\ge F_{\mathrm{cal}}(\hat q)-\kappa.
\]
Taking expectations and applying Lemma~1,
\[
\Pr\{E\le\hat q\}
\ge \mathbb E[F_{\mathrm{cal}}(\hat q)]-\kappa
\ge \frac{k}{n+1}-\kappa.
\]
The second inequality in \eqref{eq:a3-2.1} follows from
$k\ge(n+1)(1-\alpha)$; negative values are replaced by $\max(0,\cdot)$. If $k>n$,
then $\hat q=+\infty$ and the coverage is $1$. Note that the second inequality is a
\emph{marginal} statement after taking expectations: the left-hand side
$\Pr\{E\le\hat q\}$ is the constant $\gamma$, while the intermediate
$F_{\mathrm{cal}}(\hat q)$ is random. One must not read $k/(n+1)$ as a pathwise
lower bound on $F_{\mathrm{cal}}(\hat q)$.

\subsubsection{Proof of Theorem~2B}

Assume first that $k\le n$. If $k>n$, then $\hat q=+\infty$ and
$F_{\mathrm{test}}(\hat q)=1$, while the right-hand side of \eqref{eq:a3-2.2} does
not exceed $1$ (for $n\ge3$, $\varepsilon_{\mathrm{cal}}\ge1/n$ because
$\sqrt{\ln2/(2n)}\ge1/n$ iff $n\ge2/\ln2\approx2.885$; for $n=1,2$ the claim is
immediate), so the conclusion is trivial. Define the events
$B=\{\sup_x|\hat F_{\mathrm{cal}}-F_{\mathrm{cal}}|\le\varepsilon_{\mathrm{cal}}\}$
and $A=\{\kappa\le\hat\kappa+\varepsilon_\kappa\}$. By Lemma~2,
$\Pr\{B\}\ge1-\delta_{\mathrm{cal}}$; by Lemma~3, $\Pr\{A\}\ge1-\delta_\kappa$;
hence $\Pr\{A\cap B\}\ge1-\delta_{\mathrm{cal}}-\delta_\kappa$.

On $B$, $F_{\mathrm{cal}}(\hat q)\ge\hat F_{\mathrm{cal}}(\hat q)-\varepsilon_{\mathrm{cal}}$;
since $\hat q$ is the $k$-th order statistic,
$\hat F_{\mathrm{cal}}(\hat q)\ge k/n$ (this remains true with ties). On $A$,
$\kappa\le\hat\kappa+\varepsilon_\kappa$. Therefore
\[
F_{\mathrm{test}}(\hat q)
\ge F_{\mathrm{cal}}(\hat q)-\kappa
\ge \frac{k}{n}-\varepsilon_{\mathrm{cal}}-\hat\kappa-\varepsilon_\kappa.
\]
Taking probabilities over $A\cap B$ gives \eqref{eq:a3-2.2}. An unconditional
marginal bound cannot be read directly from the same display: one either takes a
further expectation as in Remark~2.3, or combines the constant bound of 2A with
$A$ as in Corollary~2.2 (the latter is stronger but applies only to the marginal
coverage rate $\gamma$).

\subsubsection{Implementation notes and limitations}

\begin{enumerate}
\item \textbf{Leave-fold-out split conformal, not CV+.} The code
  \path{_cross_tau_detail} (\path{policy_value_check.py}:1982--1987) takes, for
  each fold $f$, the calibration set \path{cal = cal_scores[cal_folds != f]} and
  the quantile \path{q = conformal_quantile(cal, alpha)}, and declares coverage
  when \path{test_scores[test_folds == f] <= q}. Thus the $\hat q$ of every test
  point \emph{is} the $k$-th order statistic of a single split, with
  $n_{\mathrm{cal}}=100$ or $104$ ($k=91$ or $95$), not a CV+ aggregate quantile.
  Theorem~2A applies exactly per fold, provided the scores within a $\tau$ are
  i.i.d.; the fold assignment groups the four windows of the same seed into one
  fold, so the calibration and test sides come from different seeds, and i.i.d.\
  within the calibration set invokes the standing assumption that the 128 paths
  are mutually independent. The five fold coverages average five valid
  split-conformal coverages. The figures with $n=m=128$ are a balanced reference;
  with the pooled $\hat\kappa$, the exact per-fold 2A bounds are $0.784$--$0.788$
  for the studentized score (versus $0.790$); with the strict per-fold
  $\hat\kappa$ they fall to $0.552$--$0.783$ (Table~\ref{tab:a3-strict}). The real
  limitations are not that ``the quantile is not an order statistic'', but that
  $\hat\kappa$ in the pooled row is pooled over all $128+128$ scores and overlaps
  the evaluated test fold, that the five folds share data and are correlated, and
  that a single dataset cannot validate the 95\% PAC frequency.
\item \textbf{Paired samples.} The scores at $\tau=2$ and $\tau=81$ are paired by
  path. The leave-fold-out protocol holds out entire paths, so the calibration set
  (scores at $\tau=2$ from the other folds) and the test points (scores at
  $\tau=81$ from the held-out fold) come from different paths and are independent;
  pairing does \emph{not} break the independence required by the theorem. It
  affects the bootstrap SE of $\hat\kappa$ (\texttt{bootstrap\_ks} resamples paths
  in pairs, which is correct) and the fact that all numbers come from the same
  paths. If $\hat\kappa$ is pooled over all 128 paths, the two one-sample DKW
  bounds in Lemma~3 still hold separately (the union bound does not require the
  two samples to be independent). The strict version uses the fold's own
  calibration set ($n_{\mathrm{cal}}$ scores at $\tau=2$) and the fold's own
  $m=24$--$28$ scores at $\tau=81$; overlap of $\hat\kappa$ with the evaluated
  test fold does \emph{not} affect the validity of \eqref{eq:a3-2.2} or
  \eqref{eq:a3-2.6} (both $F_{\mathrm{test}}(\hat q)$ and $\gamma$ are population
  quantities); non-overlap is required only for the PAC-frequency validation in
  item~7. Table~\ref{tab:a3-strict} lists the five folds. The 2B values are
  $0.345$--$0.354$ with the pooled $\hat\kappa$ and $m=128$, $0.170$--$0.184$ with
  the fold's $m$ but the pooled $\hat\kappa$ (the old ``strict'' row), and
  $0.000$--$0.165$ with the strict per-fold $\hat\kappa$; only the last is the
  strict row, and fold~3 (zero-based) is empty.
\item \textbf{Fixed model.} The score functions (the $h$ of G0 and
  $\kappa/\mu/\text{jump}$) are fixed between calibration and test. If they are
  re-estimated within folds, the coverage guarantee is conditional on the fitted
  model and does not cover model uncertainty.
\item \textbf{Marginal, conditional, PAC and expectation must not be conflated.}
  2A is an exact finite-sample marginal lower bound with known $\kappa$; Corollary~2.2 is a confidence
  lower bound on the \emph{marginal coverage rate} $\gamma$ with estimated
  $\kappa$; 2B is a PAC bound on \emph{conditional} coverage given the data; and
  Remark~2.3 is an unconditional bound in expectation. They are different objects
  and cannot substitute for one another. The statistic $\hat\kappa$ is a KS
  distance between marginal score distributions; within-group coverage
  (EXTENDED/WEEKEND etc.) is only a descriptive diagnostic, not a conditional
  guarantee. The boundary for branch-conditional coverage is in
  Theorem~3/Proposition~2.
\item \textbf{Unlabeled deployment.} The estimate $\hat\kappa$ requires test
  scores. If no immediate labels are available at deployment, one can only use a
  proxy distribution or a conservative upper bound.
\item \textbf{The supremum over $x$ is conservative.} Under one-sided drift (for
  example the scale change from $\tau=2$ to $\tau=81$), the sup-norm $\kappa$
  exceeds the local CDF difference at the threshold, so the bound is loose; this
  is the price paid for estimability.
\item \textbf{Single-dataset PAC frequency versus the independent check.} The
  original batch of 128 paths and five correlated folds does not validate the
  95\% frequency. Table~A3 therefore uses 100 independent three-sample
  replicates ($n=m=64$, final $N=256$), with the radius, $\hat\kappa$, and
  final evaluation separated by independent seeds. The below-bound column
  compares noisy final-sample coverage with the bound; it is a tightness
  diagnostic, not a direct estimate of the population event
  $\{F_{\mathrm{test}}(\hat q)\ge b\}$ or an exact PAC-violation frequency.
\end{enumerate}

\begin{table}[t]
\centering
\small
\caption{Strict per-fold quantities for the Wide rebalance studentized score
($\tau=2\to81$, $\alpha=0.1$). $\hat\kappa$ is the two-sample KS distance between
the fold's own calibration scores at $\tau=2$ and its own test scores at
$\tau=81$; 2A treats $\hat\kappa$ as known, 2B uses
$\delta_{\mathrm{cal}}=\delta_\kappa=0.025$, and Cor.~2.2 uses $\delta_\kappa=0.05$.
Source: \texttt{results/theorem2\_bounds.json}, written by
\texttt{code/theorem2\_bounds.py}.}
\label{tab:a3-strict}
\begin{tabular}{@{}rrrrrrr@{}}
\toprule
Fold & $n_{\mathrm{cal}}$ & $m$ & $k$ & $\hat\kappa$ & 2A & 2B / Cor.~2.2 \\
\midrule
0 & 100 & 28 & 91 & 0.186 & 0.715 & 0.116 / 0.288 \\
1 & 100 & 28 & 91 & 0.256 & 0.645 & 0.046 / 0.218 \\
2 & 104 & 24 & 95 & 0.122 & 0.783 & 0.165 / 0.336 \\
3 & 104 & 24 & 95 & 0.353 & 0.552 & 0.000 / 0.105 \\
4 & 104 & 24 & 95 & 0.154 & 0.751 & 0.133 / 0.304 \\
\bottomrule
\end{tabular}
\end{table}

\paragraph{Independent three-sample check.} Table~\ref{tab:14} separates radius
construction, drift estimation, and final evaluation into three independently
seeded samples. This checks out-of-sample coverage and tightness without using
final-test scores to compute $\hat\kappa$. It does not convert the finite final
sample into a direct observation of the population PAC event.
\begin{table}[!ht]
\centering
\scriptsize
\caption{Three-sample out-of-sample check of Theorem 2}
\label{tab:14}
\begin{adjustbox}{max width=\linewidth}
\begin{tabular}{@{}llllllll@{}}
\toprule
Direction & Score & Mean $\hat\kappa$ & Mean 2B & Mean Cor 2.2 & Final coverage [mean CI95] & Mean width & Below 2B \\
\midrule
2$\to$81 & raw & 0.8453 & 0.0000 & 0.0000 & 0.1046 [0.0977, 0.1116] & 0.00126 & 0/100 \\
2$\to$81 & studentized & 0.1397 & 0.1989 & 0.3980 & 0.9099 [0.9029, 0.9168] & 0.02231 & 0/100 \\
2$\to$81 & CQR & 0.6112 & 0.0000 & 0.0040 & 0.8886 [0.8831, 0.8940] & 0.01910 & 0/100 \\
81$\to$2 & raw & 0.8448 & 0.0000 & 0.0000 & 1.0000 [1.0000, 1.0000] & 0.02369 & 0/100 \\
81$\to$2 & studentized & 0.1484 & 0.1902 & 0.3892 & 0.9059 [0.8972, 0.9146] & 0.00115 & 0/100 \\
81$\to$2 & CQR & 0.6052 & 0.0000 & 0.0046 & 0.9081 [0.8840, 0.9322] & 0.00199 & 0/100 \\
\bottomrule
\end{tabular}
\end{adjustbox}
\vspace{0.5em}
\par\noindent\begin{minipage}{\linewidth}\raggedright{\scriptsize Notes. Wide identified pipeline, 100 independent outer replicates. Each replicate uses 64 independently seeded calibration paths, the same number of independent target-drift paths, and 256 final target-test paths. One path is drawn per seed and window types cycle by seed. The radius uses calibration only; $\hat\kappa$ uses calibration and the drift sample only; the final sample evaluates coverage and width without re-entering either calculation. The bracket is a 95\% interval for the mean final coverage across independent outer replicates, not a binomial interval over pooled events. \texttt{Below 2B} counts replicates whose noisy final-sample coverage falls below the 2B bound; it is an empirical tightness diagnostic, not an exact PAC-violation frequency. Bounds can be empty (zero) for any score. Width is in target-return units. Both $\tau=2\to81$ and $\tau=81\to2$ are shown.}\end{minipage}\par

\end{table}


\paragraph{Relation to the other results.}
Theorem~1 gives population identification and plug-in consistency for $V^\pi$.
Theorem~2A and Corollary~2.2 give \emph{marginal} coverage at the score level
(exact bound and confidence bound with estimated $\kappa$), and Theorem~2B gives
\emph{conditional} (PAC) coverage. The two are on different tracks: Theorem~1 does
not imply a finite-sample rate, and Theorem~2 does not rely on the identification
assumptions (only on a fixed score function). Theorem~3/Proposition~2 show that
marginal coverage cannot be upgraded to branch-conditional coverage; the KS bound
on $\kappa$ concerns only the \emph{marginal} score distribution and cannot be used
to refute those impossibility results. Theorem~4 (Proposition~11) is a
within-class non-identification result on the identification side: in the
non-Markov freeze class with a bridge reward not determined by endpoints,
dropping I5/I5$'$ yields observationally equivalent models with different
values. It does not assert general necessity of I5 and has no direct
dependence on Theorem~2.

\let\theHequation\aIIItheHequation
